\chapterhead{3}{Salt That Disappears}{Dissolving, and what the scale says}\label{ch3}

Stir a spoonful of salt into water. The white grains vanish and the water is clear again. Where did the salt go? Is it gone for good?

Salt that has spread out among the water, too small to see, is \textbf{dissolved}. The water with the dissolved salt in it is a \textbf{solution}. A solution is a mixture whose parts you cannot see or sieve apart. But the parts are still there, and a scale can tell.

\begin{tryit}{Weigh before and after}
\textbf{You need:} a kitchen scale, a clear cup, water, table salt, a spoon.
\begin{enumerate}[leftmargin=1.6em]
\item Put the cup on the scale and press the button that sets the scale to zero, or write the cup's weight. Pour in 100~g of water.
\item Put a small spoonful of salt on the scale in a saucer and weigh it. Write the weight.
\item Add the salt to the water without stirring. Weigh the cup, water and salt together.
\item Stir until the salt seems gone. Weigh again.
\end{enumerate}
\par\medskip
\renewcommand{\arraystretch}{1.6}
\begin{tabular}{|l|c|c|c|c|}
\hline
 & \textbf{Cup} & \textbf{Water} & \textbf{Salt} & \textbf{Total (g)}\\ \hline
\textbf{Before stirring} & & & & \\ \hline
\textbf{After stirring} & & & & \\ \hline
\end{tabular}
\end{tryit}
\adult{Never taste anything from an experiment, not even the salt water. Wipe up spills.}

\begin{worked}{Practice numbers, not from a real test}
Cup 40~g. Water 100~g. Salt 5~g. Before stirring: $40+100+5=145$~g.\\
After stirring the cup still read 145~g. The salt looked gone, but its weight had not gone anywhere.
\end{worked}

\begin{center}
\begin{tikzpicture}
  \node[bigcir,minimum size=2.4cm] (a) at (0,0) {}; \node[font=\sffamily\scriptsize,align=center] at (0,0) {water 100\\salt 5};
  \node[bigcir,minimum size=2.4cm] (b) at (5.4,0) {}; \node[font=\sffamily\scriptsize,align=center] at (5.4,0) {solution\\105};
  \draw[moveA] (a) -- node[above,note] {stir} (b);
  \node[note,anchor=west,text width=3.2cm,align=left] at (7.1,0) {\textbf{kept:} the total weight\\\textbf{set aside:} the look of the salt};
\end{tikzpicture}
\end{center}

\begin{multix}[hmm]\multixbot[hmm]\tcblower
The salt dissolved, so it vanished. The cup should weigh the same as 100 grams of water plus a cup. I weighed it myself and I trust my idea. Do the numbers agree?
\end{multix}

\trust{We weighed the cup, the water and the salt before stirring and the whole cup after stirring.}
{The weight did not change when the salt dissolved. If the salt had stopped existing, the cup would have got lighter.}
{Where the salt is, or how much a spoonful of salt weighs for a different spoon. We only know that its weight stayed in the cup.}

\begin{keepbox}
We looked at salt in water. \textbf{We kept} \wline[3.4cm] \textbf{and we set aside} \wline[3.4cm].
\end{keepbox}

\clearpage

\subhead{Get the salt back}
If the salt is still there, we should be able to get it back. Water leaves by evaporation (Book~5), so the salt should stay.

\begin{tryit}{Evaporate the salt water}
Weigh an empty saucer. Pour about 25~g of salt water into it. Put it on a sunny sill for a few days. When it looks dry, weigh the saucer again.\par\medskip
\renewcommand{\arraystretch}{1.6}
\begin{tabular}{|l|c|}
\hline
\textbf{Empty saucer (g)} & \\ \hline
\textbf{Saucer after the water has gone (g)} & \\ \hline
\textbf{Crust of salt (g)} & \\ \hline
\end{tabular}
\end{tryit}

\begin{worked}{Practice numbers}
Empty saucer: 30~g. After three sunny days: 33~g. Crust: $33-30=3$~g.\\
The crust was white and hard, and it was the salt left behind when the water evaporated.
\end{worked}

\subhead{How much will dissolve?}
\begin{tryit}{Spoonfuls till it stops}
Add level spoonfuls of salt to 100~g of water, stirring until each is gone. Count the spoonfuls until some stays at the bottom. How many did you add? \wline[2cm]
\end{tryit}

\begin{worked}{Practice numbers}
5 spoonfuls dissolved. The 6th stayed as grains at the bottom. The water had taken in as much salt as it could.
\end{worked}

\begin{multix}[hmm]\multixbot[hmm]\tcblower
Clear water must be pure water. Salt water looks clear, so I can tell it is plain water. Is that true?
\end{multix}

\trust{We got the salt back by evaporation and weighed the crust.}
{A solid crust of salt remained, so something solid really was dissolved in the clear water.}
{That we got back every gram. Some crust sticks to the saucer, and the scale is only exact to about 1~g.}

\clearpage

\begin{puzzle}{Total weight}
A cup weighs 40~g, the water 100~g and the salt 7~g. What does the cup weigh after the salt dissolves? \wline[1.6cm]~g
\end{puzzle}

\begin{puzzle}{Crust}
An empty saucer weighs 28~g. After the salt water dried it weighs 33~g. How much salt was left? \wline[1.6cm]~g
\end{puzzle}

\begin{puzzle}{Clear is not pure}
Two cups of clear water look the same. One has dissolved salt. Without tasting, how could you find out which? \wline[6cm]
\end{puzzle}

\begin{puzzle}{Spoonfuls}
In a practice test 5 spoonfuls dissolved and the 6th did not. Write the sentence that fits the data: \wline[6cm]
\end{puzzle}

\begin{puzzle}{Spot the unfair test}
Sam added salt to cold water and Mia added salt to warm water. Mia's salt dissolved quicker. Which two things differ? \wline[6cm]
\end{puzzle}

\clearpage

\begin{puzzle}{Draw the circles}
The solution weighs 105~g. The salt was 5~g. Fill in the circle.
\begin{center}
\begin{tikzpicture}
  \node[cir] (a) at (0,0) {105}; \node[cir,dashed] (b) at (4,0.8) {}; \node[cir] (c) at (4,-0.8) {5};
  \draw[popA] (a) -- node[above,note] {evaporate} (b); \draw[popA] (a) -- (c);
  \node[note,anchor=west] at (4.6,0.8) {water left in the cup/air}; 
\end{tikzpicture}
\end{center}
\end{puzzle}

\begin{puzzle}{Sentence check}
Mark S if the data supports it, G if it is a guess about why, and B if it is too big for the data.\\
(a) ``The cup weighed the same after stirring.'' \wline[1cm]\\
(b) ``The salt broke into tiny pieces.'' \wline[1cm]\\
(c) ``Nothing ever disappears.'' \wline[1cm]
\end{puzzle}

\begin{puzzle}{Another solid}
Name another solid that might dissolve in water, and one that might not. Say how you could find out without tasting. \wline[6cm]
\end{puzzle}

\begin{puzzle}{Short on the scale}
After the salt water dried, the crust weighed 2~g, but the salt you added weighed 3~g. Name two reasons for the missing gram. \wline[6cm]
\end{puzzle}

\begin{puzzle}{Your test}
Plan a fair test of whether stirring faster dissolves salt sooner. Change: \wline[3cm] Keep the same: \wline[4cm] Measure: \wline[3cm]
\end{puzzle}
\clearpage

\thispagestyle{fancy}
\subhead{Salt crust diary}
Keep a record of getting the salt back. Weigh the saucer each day.\par\medskip
\renewcommand{\arraystretch}{3.0}
\begin{tabularx}{\linewidth}{|>{\sffamily\small\bfseries}p{1.4cm}|X|X|X|}
\hline
Day & Saucer and contents (g) & What it looks like & Crust so far (g) \\ \hline
0 & & & \\ \hline
1 & & & \\ \hline
2 & & & \\ \hline
3 & & & \\ \hline
4 & & & \\ \hline
5 & & & \\ \hline
\end{tabularx}

\vspace{8pt}
When did the weight stop dropping? \wline[6cm]\\[10pt]
Why does it stop? \wline[7cm]\\[10pt]
Draw the saucer on the first and last day.\par\medskip
\begin{tikzpicture}\draw[ink!70,line width=1pt,rounded corners=3pt] (0,0) rectangle (5.5,4.6); \draw[ink!70,line width=1pt,rounded corners=3pt] (5.9,0) rectangle (11.4,4.6);\end{tikzpicture}
\clearpage
